\documentclass[12pt, a4paper]{ctexart}
\usepackage[margin=2cm]{geometry}
\usepackage{libertine}

\usepackage{titlesec}
\usepackage{zhnumber}
\titleformat*{\section}{\Large\bfseries\raggedright}
\renewcommand{\thesection}{\zhnum{section}}
\renewcommand{\thesubsection}{\arabic{subsection}}

\setcounter{secnumdepth}{4}
\renewcommand{\theparagraph}{\thesubsubsection.\arabic{paragraph}}
\renewcommand{\paragraph}[1]{
    \refstepcounter{paragraph}
    {\bfseries\theparagraph\quad#1}\par
    \vspace{2pt}
    \noindent
}

\usepackage{enumitem}
\setlist[enumerate]{itemsep=2pt, parsep=0pt, partopsep=0pt, topsep=2pt}
\setlist[itemize]{itemsep=2pt, parsep=0pt, partopsep=0pt, topsep=2pt}
\linespread{1.2}

\usepackage{graphicx}

\usepackage{minted}
\setminted{
    breaklines=true, escapeinside=||, fontsize=\small, frame=lines,
    mathescape=$$, numbers=none, style=bw, tabsize=4
}

\usepackage{siunitx}

\begin{document}
    
\pagestyle{plain}
\thispagestyle{empty}

\noindent
\begin{tabular*}{\textwidth}{l @{\extracolsep{\fill}} r @{\extracolsep{6pt}} l}
    \LARGE{\textbf{实验报告}} & 计算机网络 & \textit{Computer Networking} \\
\end{tabular*}\\\\
\begin{tabular*}{\textwidth}{l l}
    \textbf{报告标题: } & \textbf{交换机部署及配置} \\
    \textbf{学号: } & 19240212 \\
    \textbf{姓名: } & 华博文 \\
    \textbf{日期: } & 2025 年 11 月 24 日 \\
\end{tabular*}\\
\rule[2ex]{\textwidth}{2pt}

\section{实验目的}

本实验旨在掌握 IPMininet 仿真网络的启动方法，熟悉 Open vSwitch 的基础操作；掌握主机IP地址的配置流程，理解交换机模式切换（Secure 模式与普通模式）对数据转发的影响；深入理解交换机 CAM 表的工作原理，包括 MAC 地址与端口的映射学习过程及数据流转发逻辑；通过网络连通性测试与CAM表分析，强化对二层交换网络通信机制的认知。

\section{实验内容简要描述}

\begin{enumerate}
    \item \textbf{启动 IPMininet 仿真网络}：启动 Open vSwitch 服务，通过 \mintinline{bash}`mn` 命令创建包含 1 台交换机、2 台主机的基础仿真网络；
    
    \item \textbf{配置主机IP地址}：打开主机终端，查看初始 IP 配置，修改主机网卡的 IP 地址与子网掩码至实验要求网段；
    
    \item \textbf{配置交换机}：打开交换机终端，查看交换机模式，退出 Secure 模式以启用普通二层交换功能；
    
    \item \textbf{分析CAM表及数据流}：通过 \mintinline{text}`ping` 命令测试主机间连通性，查看交换机 CAM 表，分析 MAC 地址与端口的映射关系及数据流转发过程。
\end{enumerate}

\section{实验步骤与结果分析}

\subsection{启动 IPMininet 仿真网络}

右键打开终端，执行以下命令启动 Open vSwitch 服务（指定 pid 文件并后台运行）：

\begin{minted}{bash}
ovs-vswitchd --pidfile --detach
\end{minted}

在终端中执行以下命令，启动包含 1 台交换机（s1）、2 台主机（h1、h2）的 IPMininet 仿真网络：

\begin{minted}{bash}
sudo mn --controller=remote
\end{minted}

终端输出显示网络创建过程（添加控制器、主机、交换机、链路），说明网络创建完成。

\begin{figure}[H]
    \centering
    \includegraphics[width=\textwidth]{./img/screenshot_0000.png}
\end{figure}

\subsection{配置主机 IP 地址}

打开主机 h1 的终端：

\begin{minted}{bash}
xterm h1
\end{minted}

在 h1 的终端中查看初始 IP 配置：

\begin{minted}{bash}
sudo ifconfig
\end{minted}

输出显示 h1 的 \mintinline{text}`h1-eth0` 网卡初始 IP 为 \mintinline{text}`10.0.0.1`，子网掩码为 \mintinline{text}`255.0.0.0`，不符合实验要求。

在 h1 的终端中，执行以下命令配置 IP 与子网掩码：

\begin{minted}{bash}
sudo ifconfig h1-eth0 10.0.0.2
sudo ifconfig h1-eth0 netmask 255.255.255.0
\end{minted}

再次执行 \mintinline{text}`sudo ifconfig` 验证，确认 \mintinline{text}`h1-eth0` 的 IP 已更新为 \mintinline{text}`10.0.0.2/24`。

\begin{figure}[H]
    \centering
    \includegraphics[width=\textwidth]{./img/screenshot_0001.png}
\end{figure}

打开 h2 的终端：

\begin{minted}{bash}
xterm h2
\end{minted}

在 h2 的终端中执行以下命令，配置 IP 为 \mintinline{text}`10.0.0.3/24`：

\begin{minted}{bash}
sudo ifconfig h2-eth0 10.0.0.3
sudo ifconfig h2-eth0 netmask 255.255.255.0
\end{minted}

执行 \mintinline{text}`sudo ifconfig` 验证，确认配置符合实验要求。

\begin{figure}[H]
    \centering
    \includegraphics[width=\textwidth]{./img/screenshot_0002.png}
\end{figure}

\subsection{配置交换机}

打开交换机 s1 的终端：

\begin{minted}{bash}
xterm s1
\end{minted}

在 s1 的终端中查看交换机信息：

\begin{minted}{bash}
ovs-vsctl show
\end{minted}

输出显示交换机的 \mintinline{text}`fail_mode` 为 \mintinline{text}`secure`，说明需依赖 SDN 控制器才能转发数据。

在 s1 的终端中，执行以下命令退出 Secure 模式：

\begin{minted}{bash}
ovs-vsctl del-fail-mode s1
\end{minted}

再次执行 \mintinline{text}`ovs-vsctl show` 验证，\mintinline{text}`fail_mode` 字段消失，说明交换机已切换为普通二层交换模式。

\begin{figure}[H]
    \centering
    \includegraphics[width=\textwidth]{./img/screenshot_0003.png}
\end{figure}

\subsection{分析 CAM 表及数据流}

在 h1 的终端中，测试与 h2 的连通性（发送 4 个 ICMP 包）：

\begin{minted}{bash}
ping 10.0.0.3 -c 4
\end{minted}

输出显示 4 个包均成功接收，丢包率为 0\%，说明 h1 与 h2 连通正常。

在 h2 的终端中，测试与 h1 的连通性：

\begin{minted}{bash}
ping 10.0.0.2 -c 4
\end{minted}

输出显示连通正常，丢包率为 0\%。

\begin{figure}[H]
    \centering
    \includegraphics[width=\textwidth]{./img/screenshot_0004.png}
\end{figure}

在 s1 的终端中，执行以下命令查看 CAM 表（MAC 地址表）：

\begin{minted}{bash}
ovs-appctl fdb/show s1
\end{minted}

输出显示 CAM 表包含两条记录：端口 1 对应 h1 的 MAC 地址，端口 2 对应 h2 的MAC地址，说明交换机已学习到主机 MAC 与端口的映射关系。

\begin{figure}[H]
    \centering
    \includegraphics[width=\textwidth]{./img/screenshot_0005.png}
\end{figure}

\section{实验中遇到的问题及体会}

在启动 ovs-vswitchd 服务时，因未使用 sudo 权限导致命令执行失败，意识到 Open vSwitch 服务需要管理员权限；配置主机 IP 时，误将网卡名称写为 ``eth0'' 而非 mininet 中的 ``h1-eth0''，导致配置不生效，这让我注意到 mininet 中主机网卡的命名规则。

通过本次实验，我直观认知了二层交换网络的工作机制：交换机通过学习数据包源 MAC 构建 CAM 表，再根据目的 MAC 查询表项转发数据，实现精准转发而非广播；Secure 模式体现了 SDN 的集中管控特性，普通模式则回归传统交换机的自主转发逻辑。

IPMininet 为网络实验提供了低成本、可重复的环境，通过终端操作，我将 ``MAC 地址学习'' ``数据转发'' 等理论与实践结合，深化了对二层网络的理解。后续实验需更关注命令细节与组件命名规则，提升配置准确性与故障排查效率。

\end{document}